\section{Weights for exceptional sets}
\label{sec:weights}

The next proposition supplies a cover at a fixed fraction of the original
density. Its proof combines an orthogonal expansion with a signed measure
on arrays of sets. The expansion bounds the total weight, while the signed
measure forces a large sum of high powers on every exceptional set.
Lemma~\ref{lem:covering} then converts these two bounds into a cover.

\begin{proposition}\label{prop:half-density}
Let $L=2^{70}$. For every integer $N\ge1$, every $q\in(0,1)$, and every
$\mathcal F\subseteq 2^{[N]}$ satisfying $\mu_q(\mathcal F)\ge1/2$,
the family $E_{32}(\mathcal F)$ is $(q/L)$-small.
\end{proposition}

\begin{proof}
Identify subsets of $[N]$ with their indicator vectors and write
$f=\mathbf 1_{\mathcal F}$. For $U\subseteq[N]$, define
\begin{equation}\label{eq:weights-definition}
 b(U)=(1-q)^{|U|}
 \sum_{z\in\{0,1\}^{U}}(-1)^{\sum_{i\in U}z_i}
 \mathbb E[f(z,Y)],
 \qquad w(U)=b(U)^2,
\end{equation}
where $Y$ has independent Bernoulli-$q$ coordinates on $[N]\setminus U$.
In particular, $b(\varnothing)=\mu_q(\mathcal F)$.

We will compare a global bound on the weights with a lower bound on
their sixteenth powers over subsets of each exceptional set.
In a bin where $w(U)\le2^{-h}$, fewer than $2^{12h}$ terms contribute
at most $2^{-4h}$ to the sum of sixteenth powers. This is the gap
exploited by Lemma~\ref{lem:covering}.

\paragraph{The global weight bound.}
We use the standard biased Fourier basis
\cite[Section~8.4]{ODonnell14}, with its normalization made explicit here.
Under the product law $\mu_q$, the functions
\[
 \chi_U(x)=\frac{\prod_{i\in U}(q-x_i)}{[q(1-q)]^{|U|/2}},
 \qquad U\subseteq[N],
\]
are orthonormal: each standardized coordinate has mean zero
and variance one, and the coordinates are independent. Since $0<q<1$, the measure has
full support; these $2^N$ orthonormal functions form a basis of the
$2^N$-dimensional space of real functions on the cube. Set
$\widehat f(U)=\mathbb E[f(X)\chi_U(X)]$ for $X\sim\mu_q$.
For a Bernoulli-$q$ bit $X_i$ and $z_i\in\{0,1\}$,
\[
 \mathbb P(X_i=z_i)(q-z_i)=q(1-q)(-1)^{z_i}.
\]
Summing over the coordinates in $U$ therefore gives
\[
 \widehat f(U)
 =[q(1-q)]^{|U|/2}
 \sum_{z\in\{0,1\}^{U}}(-1)^{\sum_{i\in U}z_i}
 \mathbb E[f(z,Y)]
 =\left(\frac{q}{1-q}\right)^{|U|/2}b(U).
\]
Consequently, Parseval's identity yields
\begin{equation}\label{eq:global-weight-bound}
 \begin{aligned}
 \sum_{U\subseteq[N]}q^{|U|}w(U)
 &=\sum_{U\subseteq[N]}(1-q)^{|U|}\widehat f(U)^2\\
 &\le\sum_{U\subseteq[N]}\widehat f(U)^2
 =\mathbb E[f(X)^2]
 =\mu_q(\mathcal F)\le1.
 \end{aligned}
\end{equation}
This calculation holds for all $q\in(0,1)$.

\paragraph{A bound on every exceptional set.}
Fix $S\in E_{32}(\mathcal F)$. We claim that
\begin{equation}\label{eq:exceptional-weight-bound}
 \sum_{\varnothing\ne U\subseteq S}w(U)^{16}\ge2^{-32}.
\end{equation}
Consider an array of bits with $32$ rows and $N$ columns. For a column
$y\in\{0,1\}^{32}$, let
\[
 P(y)=\prod_{j=1}^{32}q^{y_j}(1-q)^{1-y_j}.
\]
For columns in $S$, we want a correction that cancels the all-zero column
and factors across rows. A suitably scaled parity term has both properties.
Give each column outside $S$ the mass function $P$, and give each column
in $S$ the signed mass function
\[
 R(y)=P(y)-(1-q)^{32}(-1)^{\sum_{j=1}^{32}y_j}.
\]
With two rows, this is precisely Li's signed kernel after multiplying
by the Bernoulli masses~\cite[Section~2]{Li26}. Park also notes the
extension of Li's argument to an even number of unions
\cite[Corollary~3.5, footnote~3]{Park26}. Here the higher powers obtained
from 32 rows are combined with Lemma~\ref{lem:covering} to obtain a
fixed density loss.
The signed-operator approach to forbidden configurations and its biased
Fourier basis have earlier forms in
\cite[Section~2.2]{Friedgut08} and
\cite[Section~4.1]{EllisFilmusFriedgut12}.
Let $\nu_S$ be the product of these column masses. Since
$R(0,\ldots,0)=0$, every array with nonzero $\nu_S$-weight has a $1$
in each column of $S$. The sets represented by its rows therefore have
union containing $S$. By the definition of $E_{32}(\mathcal F)$, these
rows cannot all lie in $\mathcal F$. Writing $x_j$ for row $j$, we obtain
\[
 \sum_{x\in\{0,1\}^{32\times N}}\nu_S(x)
 \prod_{j=1}^{32}f(x_j)=0.
\]
This uses pointwise vanishing of the product on the support of $\nu_S$;
nonnegativity of the signed masses is unnecessary.

Expand the column products, and let $U\subseteq S$ be the columns in
which the subtracted term is chosen. After taking out $(-1)^{|U|}$,
the correction contributes
\[
 (1-q)^{32|U|}\prod_{j=1}^{32}
 (-1)^{\sum_{i\in U}x_{ji}}.
\]
All coordinates outside $U$ retain their Bernoulli product masses.
Thus the sum factors over rows. In each row the factor is precisely
$b(U)$ from \eqref{eq:weights-definition}, because the power
$(1-q)^{32|U|}$ distributes as $(1-q)^{|U|}$ to each row.
It follows that
\begin{equation}\label{eq:signed-identity}
 0=\sum_{U\subseteq S}(-1)^{|U|}b(U)^{32}.
\end{equation}
All products and sums here are finite. Isolating the empty term and
using the triangle inequality gives
\[
 2^{-32}\le b(\varnothing)^{32}
 \le\sum_{\varnothing\ne U\subseteq S}|b(U)|^{32}
 =\sum_{\varnothing\ne U\subseteq S}w(U)^{16},
\]
which proves \eqref{eq:exceptional-weight-bound}.

\paragraph{Constructing the cover.}
Set $\rho=q/L$. First include all sets in
\[
 \mathcal G_0=
 \{U\subseteq[N]:U\ne\varnothing,\ w(U)\ge2^{-64|U|}\}
\]
as generators. For $U\in\mathcal G_0$ with $r=|U|\ge1$,
\[
 \rho^r=q^r2^{-70r}
 \le2^{-6r}q^r w(U)
 \le2^{-6}q^r w(U).
\]
Hence \eqref{eq:global-weight-bound} gives
\begin{equation}\label{eq:outright-cost}
 c_\rho(\mathcal G_0)\le2^{-6}.
\end{equation}

Partition the remaining nonempty sets of positive weight into families
\[
 \mathcal H_{r,h}=
 \bigl\{U\subseteq[N]:|U|=r,\ U\notin\mathcal G_0,
 \quad 2^{-(h+1)}<w(U)\le2^{-h}\bigr\},
\]
where $r\ge1$ and $h$ is an integer. Every such set belongs to exactly
one bin, and every occupied bin has $h\ge64r$. Indeed, its weights are
strictly below $2^{-64r}$; a weight equal to that threshold was included
in $\mathcal G_0$. By \eqref{eq:global-weight-bound},
\[
 |\mathcal H_{r,h}|q^r2^{-(h+1)}
 <\sum_{U\in\mathcal H_{r,h}}q^r w(U)\le1
\]
for each occupied bin. In particular,
\[
 |\mathcal H_{r,h}|\rho^r\le2^{h+1}.
\]
Since $h\ge64r\ge64$, all the hypotheses of
Lemma~\ref{lem:covering} hold. Choose a cover $\mathcal G_{r,h}$ of
the sets containing at least $2^{12h}$ members of $\mathcal H_{r,h}$,
with
\[
 c_\rho(\mathcal G_{r,h})\le3\cdot2^{-h}.
\]
There are only finitely many occupied bins, because the ground set is
finite. We may therefore choose these covers simultaneously and form
\[
 \mathcal G=\mathcal G_0\cup\bigcup_{\mathcal H_{r,h}\ne\varnothing}
 \mathcal G_{r,h}.
\]
Removing any repeated generators can only decrease its cost. The
geometric sum
\begin{equation}\label{eq:bin-cost-sum}
 3\sum_{r\ge1}\sum_{h\ge64r}2^{-h}
 =\frac{6}{2^{64}-1}<\frac14
\end{equation}
and \eqref{eq:outright-cost} imply
\[
 c_\rho(\mathcal G)
 \le\frac1{64}+\frac{6}{2^{64}-1}<\frac12.
\]

It remains to show that $\mathcal G$ covers $E_{32}(\mathcal F)$.
Suppose that $S\subseteq[N]$ contains no member of $\mathcal G$.
It contains no member of $\mathcal G_0$ and fewer than $2^{12h}$
members of each bin $\mathcal H_{r,h}$. Zero weights contribute nothing,
so
\begin{equation}\label{eq:bin-power-sum}
 \begin{aligned}
 \sum_{\varnothing\ne U\subseteq S}w(U)^{16}
 &\le\sum_{r\ge1}\sum_{h\ge64r}2^{12h}2^{-16h}\\
 &=\frac{16}{15(2^{256}-1)}<2^{-32}.
 \end{aligned}
\end{equation}
By \eqref{eq:exceptional-weight-bound}, such an $S$ cannot belong to
$E_{32}(\mathcal F)$. Thus $\mathcal G$ is the required cover.
\end{proof}
